\section{ISA and simulator}

The instruction set has a two level memory model: DRAM holds inputs, constants,
and outputs; the scratchpad, a small fixed size fast memory with a default of one
megabyte, holds the working set. The opcodes are NOP, HALT, DMA\_LOAD, DMA\_STORE,
CONV2D, MATMUL, RELU, ADD, MUL, POOL\_MAX, POOL\_AVG, and RESHAPE. The binary
format is a fixed header followed by tagged, length prefixed records, which is
simpler to get right and to disassemble than a bit packed encoding while remaining
a real, documented format. The \texttt{npu-objdump} tool renders the DRAM layout
and the instruction stream.

The simulator executes the stream against DRAM and scratchpad arrays with fp32
kernels for each instruction. Its cost model is analytical, not cycle accurate.
DMA costs bytes over an assumed bandwidth; convolution and matmul cost their
multiply accumulate count over an assumed systolic throughput of 256 operations
per cycle, the arithmetic of a sixteen by sixteen array; elementwise instructions
cost their element count over a lane width; and every instruction charges a fixed
issue overhead. These constants are documented assumptions, and the cycle and DRAM
figures the simulator reports are simulated estimates.
